\documentclass[11pt,a4paper]{article}

% --- Variables du TP
\newcommand{\TPModule}{Datawarehouse}
\newcommand{\TPNumero}{3}
\newcommand{\TPKeywords}{datawarehouse, industrialisation, IA, sécurité, API, orchestration}
\newcommand{\TPSujet}{Exploitation avancée du DWH (IA)}
\newcommand{\TPTheme}{Clinique}
\usepackage[utf8]{inputenc}
\usepackage[french]{babel}
\usepackage[T1]{fontenc}
\usepackage{amsmath}
\usepackage{amsfonts}
\usepackage{xcolor}
\usepackage{amssymb}
\usepackage{graphicx}
\usepackage[left=2cm,right=2cm,top=2cm,bottom=2cm]{geometry}
\usepackage{fancyhdr}
\usepackage[bottom]{footmisc}
\usepackage{tikz}
\usepackage{fourier}
\usepackage{pst-text}
\usepackage{multicol}
\usepackage{comment}
\usepackage{array}
\usepackage{enumitem}
\setlist[itemize]{label=\textbullet} % pour ne plus utiliser les bulletpoints après \item
\setlength{\parindent}{0pt} % pour supprimer les \noindent
\usepackage{amssymb} % pour les symboles
\usepackage{tabularx}
\usepackage{tcolorbox}
\usepackage{fontawesome}
\usepackage{listings}

% --- Réglages listings (inchangé)
\lstset{
  basicstyle=\ttfamily\small,
  backgroundcolor=\color{gray!10},
  frame=single,
  rulecolor=\color{gray!40},
  breaklines=true,
  columns=fullflexible,
  showstringspaces=false,
  aboveskip=0.5em,
  belowskip=0.5em,
  keywordstyle=\color{blue},
  commentstyle=\color{gray},
}
\lstdefinelanguage{markdown}{
  sensitive=true,
  morecomment=[l]{<!--},
  morecomment=[s]{<!--}{-->},
  morekeywords={},
}

% --- Réglages listings (inchangé)
\lstset{
  basicstyle=\ttfamily\small,
  backgroundcolor=\color{gray!10},
  frame=single,
  rulecolor=\color{gray!40},
  breaklines=true,
  columns=fullflexible,
  showstringspaces=false,
  aboveskip=0.5em,
  belowskip=0.5em,
  keywordstyle=\color{blue},
  commentstyle=\color{gray},
}
\lstdefinelanguage{markdown}{
  sensitive=true,
  morecomment=[l]{<!--},
  morecomment=[s]{<!--}{-->},
  morekeywords={},
}

% --- Fancyhdr construit depuis les variables
\pagestyle{fancy}
\setlength{\headheight}{26pt}
\renewcommand{\baselinestretch}{1.3}

\fancyhead[L]{\textbf{EPSI\\\TPModule}}
\fancyhead[R]{\textbf{TP n°\TPNumero\\\TPSujet}}
\renewcommand{\arraystretch}{1}

% --- Hyperref + metadata construits depuis les variables
\usepackage[
  colorlinks=true,
  linkcolor=blue!70!black,
  urlcolor=blue!70!black,
  citecolor=blue!70!black
]{hyperref}

\hypersetup{
  pdftitle={TP n°\TPNumero\space - \TPSujet},
  pdfauthor={Thibault CLEMENT},
  pdfsubject={TP n°\TPNumero\space - \TPSujet},
  pdfkeywords={\TPKeywords},
}

% --- Composants réutilisables
\newcommand{\TPFrontPage}{%
  ~
  \begin{center}
    \framebox{\parbox[c]{11cm}{%
      \begin{center}
        \Large{\textbf{TP n°\TPNumero~:~\TPSujet\\\large{\TPTheme}}}
      \end{center}
    }}
  \end{center}
  \vspace{0.5cm}
}

% --- Tes macros custom existantes
\newcommand{\file}[1]{~\texttt{\textcolor{blue!70!black}{#1}}~}

\newtcolorbox{reflectionbox}[1][]{%
  colback=blue!5,
  colframe=blue!60,
  title={\faLightbulbO~~Réflexions},
  fonttitle=\bfseries,
  left=1mm,
  right=1mm,
  top=1mm,
  bottom=1mm,
  #1
}


\newcommand{\checkbox}{\(\square\)~~}
\graphicspath{{images/}}
\setlength{\columnseprule}{1pt}
\def\columnseprulecolor{\color{black}}



\begin{document}
\TPFrontPage

% ---------------------------------------------------------------------------------------------

\section*{Contexte}

Lors des TP précédents, vous avez :

\begin{itemize}
  \item construit un pipeline de données (RAW → BRONZE → SILVER),
  \item modélisé un datamart décisionnel (couche GOLD),
  \item réalisé un premier tableau de bord répondant aux questions du directeur.
\end{itemize}

La clinique souhaite désormais aller plus loin et professionnaliser son système.

\section*{Architecture existante}

Base de données : \file{dwh/clinic\_dwh.db}

Couches :

\begin{itemize}
  \item \file{bronze\_*}
  \item \file{silver\_*}
  \item \file{gold\_*}
\end{itemize}

Votre travail doit s'appuyer exclusivement sur la couche GOLD validée au TP2.

% =============================================================================================

\section*{Objectif}

Formaliser l’architecture actuelle et proposer une vision cible.

\subsection*{Travail demandé}

\begin{enumerate}
  \item Réaliser un schéma d’architecture globale du système :
  \begin{itemize}
    \item sources,
    \item ingestion,
    \item dbt,
    \item base SQLite,
    \item outil BI.
  \end{itemize}
  \item Identifier 3 limites actuelles du système.
  \item Proposer 3 améliorations possibles.
\end{enumerate}

\subsection*{Livrable}

\begin{itemize}
  \item Diagramme d’architecture (draw.io ou équivalent).
  \item Analyse rédigée (10 lignes).
\end{itemize}

% =============================================================================================
\section*{Parcours IA / DATA}

\subsection*{Objectif}

Proposer un prototype permettant d’anticiper les semaines en tension.

\subsection*{Travail demandé}

\textbf{1. Construction du dataset}

\begin{itemize}
  \item Définir une cible :
  \begin{itemize}
    \item \file{is\_tension\_t+1}
    \item ou \file{patients\_refused\_t+1}
  \end{itemize}
  \item Construire au moins 5 variables explicatives :
  \begin{itemize}
    \item moyenne mobile,
    \item variation semaine précédente,
    \item ratio demande / capacité,
    \item indicateur de saisonnalité,
    \item présence du personnel.
  \end{itemize}
  \item Justifier l’absence de fuite de données.
\end{itemize}

\textbf{2. Modélisation}

\begin{itemize}
  \item Séparation temporelle train/test.
  \item Implémenter un modèle simple (logistique ou RandomForest).
  \item Évaluer les performances (accuracy, recall, matrice de confusion).
\end{itemize}

\textbf{3. Interprétation}

\begin{itemize}
  \item Identifier les variables les plus influentes.
  \item Discuter des limites liées à la volumétrie.
\end{itemize}

\subsection*{Livrables}

\begin{itemize}
  \item Notebook commenté.
  \item Dataset généré.
  \item Analyse critique.
\end{itemize}

\end{document}